\chapter{Numerical Series}\label{ch:b1:series}

Summing infinitely many numbers means taking the limit of the partial
sums --- nothing more, nothing less. This chapter sets up the
definitions and the convergence tests usable in first year:
comparison and equivalents for positive terms, the ratio test, the
integral comparison giving the Riemann series, absolute convergence,
and the alternating series theorem. The finer theory (products of
series, summation by packets, series of functions) belongs to the
second year.

\section{Generalities}

\begin{definition}\label{def:b1:series:def}
Given a sequence $(u_n)$, the \emph{series}\index{series} $\sum u_n$
is the sequence of \emph{partial sums} $S_N = \sum_{n=0}^{N} u_n$.
The series \emph{converges} when $(S_N)$ converges; the limit is the
\emph{sum} $\sum_{n=0}^{\infty} u_n$, and $R_N = \sum_{n > N} u_n =
S - S_N$ is the \emph{remainder}, which tends to $0$.
\end{definition}

\begin{example}[Geometric series]\label{ex:b1:series:geometric}
For $q \in \C$: $\;S_N = \sum_{n=0}^{N} q^n = \frac{1 -
q^{N+1}}{1-q}$ ($q \neq 1$). The series converges iff $\abs q < 1$
(\cref{exo:b1:seq:3}), with\index{geometric series}
\[
\sum_{n=0}^{\infty} q^n = \frac{1}{1 - q} .
\]
\end{example}

\begin{proposition}[First facts]\label{prop:b1:series:first}
\begin{enumerate}
  \item If $\sum u_n$ converges, then $u_n \to 0$. (The converse is
        \emph{false}: the harmonic series.)
  \item Linearity: convergent series add and scale, with the expected
        sums.
  \item (Telescoping\index{telescoping series}) $\sum (v_{n+1} -
        v_n)$ converges iff $(v_n)$ converges, with sum $\lim v_n -
        v_0$.
  \item Changing finitely many terms does not affect convergence
        (only the sum).
\end{enumerate}
\end{proposition}

\begin{proof}
(1) $u_N = S_N - S_{N-1} \to S - S = 0$. The harmonic series has
$u_n = \frac1n \to 0$ yet diverges (\cref{exo:b1:seq:5}). (2)
Operations on limits. (3) $S_N = v_{N+1} - v_0$. (4) The partial
sums change by an eventually constant amount.
\end{proof}

\section{Series with nonnegative terms}

\begin{theorem}[Bounded partial sums]\label{thm:b1:series:positive}
If $u_n \geq 0$ for all $n$, the partial sums increase, so: $\sum
u_n$ converges $\iff$ its partial sums are bounded above. Hence the
\emph{comparison test}: if $0 \leq u_n \leq v_n$ for all (large)
$n$,
\[
\sum v_n \text{ converges} \implies \sum u_n \text{ converges},
\qquad
\sum u_n \text{ diverges} \implies \sum v_n \text{ diverges}.
\]
And the \emph{equivalents test}: if $u_n \sim v_n$ with $v_n \geq
0$, the two series have the same nature.
\end{theorem}

\begin{proof}
Monotone limit theorem (\cref{thm:b1:seq:monotone}) for the first
point; comparison of partial sums for the second. Equivalents: for
large $n$, $\frac12 v_n \leq u_n \leq 2 v_n$ (definition of $\sim$
with $\varepsilon = \frac12$), and comparison applies both ways.
\end{proof}

\begin{theorem}[Integral comparison; Riemann series]\label{thm:b1:series:riemann}
Let $f$ be continuous, nonnegative and \emph{decreasing} on
$\intco{1}{+\infty}$. Then
\[
\int_1^{N+1} f(t)\,\dd t \;\leq\; \sum_{n=1}^{N} f(n) \;\leq\; f(1)
+ \int_1^{N} f(t)\,\dd t ,
\]
so $\sum f(n)$ converges iff $\bigl(\int_1^x f\bigr)$ is bounded. In
particular, for $\alpha \in \R$:\index{Riemann series}
\[
\sum_{n \geq 1} \frac{1}{n^\alpha} \text{ converges}
\iff \alpha > 1,
\]
and $\sum_{n=1}^{N} \frac1n = \ln N + O(1)$.
\end{theorem}

\begin{proof}
For $n \leq t \leq n+1$: $f(n+1) \leq f(t) \leq f(n)$; integrating
over $\intcc{n}{n+1}$ and summing gives the two framings. For $f(t)
= t^{-\alpha}$ ($\alpha \neq 1$): $\int_1^x t^{-\alpha}\dd t =
\frac{x^{1-\alpha} - 1}{1 - \alpha}$, bounded iff $\alpha > 1$; for
$\alpha = 1$ the integral is $\ln x \to \infty$, and the framing
gives $\ln(N+1) \leq H_N \leq 1 + \ln N$. For $\alpha \leq 0$ the
terms do not tend to $0$.
\end{proof}

\begin{theorem}[Ratio test (d'Alembert)]\label{thm:b1:series:ratio}
Let $u_n > 0$ with $\frac{u_{n+1}}{u_n} \to \ell$.\index{ratio test}
\begin{itemize}
  \item If $\ell < 1$: $\sum u_n$ converges;
  \item if $\ell > 1$: $u_n \to +\infty$, divergence;
  \item if $\ell = 1$: no conclusion ($\sum \frac1n$ diverges, $\sum
        \frac{1}{n^2}$ converges).
\end{itemize}
\end{theorem}

\begin{proof}
If $\ell < 1$, fix $q \in \intoo{\ell}{1}$: beyond some $N$,
$u_{n+1} \leq q\,u_n$, so $u_n \leq u_N q^{\,n-N}$ by induction:
comparison with a geometric series. If $\ell > 1$: beyond some $N$
the sequence $(u_n)$ is increasing, so it cannot tend to $0$
(its limit, if any, is $\geq u_N > 0$); by
\cref{prop:b1:series:first}~(1), divergence --- and in fact $u_n
\geq u_N q^{n-N}$ with $q > 1$ gives $u_n \to \infty$.
\end{proof}

\begin{example}\label{ex:b1:series:ratio}
$\sum \frac{x^n}{n!}$ converges for every $x > 0$: ratio
$\frac{x}{n+1} \to 0$. Its sum is $\eu^x$: by Taylor--Lagrange
(\cref{thm:b1:taylor:lagrange}) on $\intcc{0}{x}$,
\[
\Bigl| \eu^x - \sum_{k=0}^{n} \frac{x^k}{k!} \Bigr|
\leq \eu^{x}\, \frac{x^{n+1}}{(n+1)!} \xrightarrow[n\to\infty]{} 0 ,
\]
the bound tending to $0$ because the factorial dominates
(\cref{exo:b1:integration:9}~(1) used the same fact). The same
argument sums the $\sin$, $\cos$, $\sinh$, $\cosh$ series on all of
$\R$.
\end{example}

\section{Absolute convergence; alternating series}

\begin{theorem}[Absolute convergence]\label{thm:b1:series:absolute}
If $\sum \abs{u_n}$ converges (\emph{absolute
convergence}\index{absolute convergence}), then $\sum u_n$ converges,
and $\bigl|\sum u_n\bigr| \leq \sum \abs{u_n}$. This holds for real
or complex terms.
\end{theorem}

\begin{proof}
The partial sums satisfy, for $M > N$ (Cauchy criterion,
\cref{thm:b1:seq:complete}):
\[
\abs{S_M - S_N} = \Bigl| \sum_{n=N+1}^{M} u_n \Bigr|
\leq \sum_{n=N+1}^{M} \abs{u_n},
\]
which is small for large $N$ since the partial sums of $\sum\abs{u_n}$
form a Cauchy sequence. So $(S_N)$ is Cauchy, hence convergent. The
inequality passes to the limit from the finite triangle inequality.
\end{proof}

\begin{theorem}[Alternating series test]\label{thm:b1:series:alternating}
Let $(a_n)$ be decreasing with $a_n \to 0$. Then the alternating
series $\sum (-1)^n a_n$ converges; its sum lies between any two
consecutive partial sums, and\index{alternating series}
\[
\abs{R_N} = \Bigl| \sum_{n > N} (-1)^n a_n \Bigr| \leq a_{N+1} .
\]
\end{theorem}

\begin{proof}
The even and odd partial sums are adjacent: $S_{2p+2} - S_{2p} =
a_{2p+2} - a_{2p+1} \leq 0$ (decreasing), $S_{2p+1} - S_{2p-1} =
a_{2p} - a_{2p+1} \geq 0$ (increasing), and $S_{2p} - S_{2p+1} =
a_{2p+1} \to 0$. By \cref{thm:b1:seq:adjacent} they share a limit
$S$, which the two subsequences criterion
(\cref{prop:b1:seq:subsequences}) makes the limit of $(S_N)$;
moreover $S$ is trapped between consecutive partial sums, and
$\abs{S - S_N}$ is at most the gap to the next one, $a_{N+1}$.
\end{proof}

\begin{example}[Alternating harmonic series]\label{ex:b1:series:ln2}
$\sum_{n \geq 1} \frac{(-1)^{n-1}}{n}$ converges (alternating test)
but not absolutely (harmonic series). Its sum is $\ln 2$: from the
finite geometric identity $\frac{1}{1+t} = \sum_{k=0}^{n-1} (-t)^k +
\frac{(-t)^n}{1+t}$, integrate over $\intcc{0}{1}$:
\[
\ln 2 = \sum_{k=1}^{n} \frac{(-1)^{k-1}}{k}
+ (-1)^n \int_0^1 \frac{t^n}{1+t}\,\dd t ,
\qquad
0 \leq \int_0^1 \frac{t^n}{1+t}\,\dd t \leq \frac{1}{n+1} \to 0 .
\]
The convergence is painfully slow ($R_N \approx \frac{1}{N}$) ---
alternating series converge by cancellation, not by smallness.
\end{example}

\begin{method}[Deciding the nature of a series]\label{met:b1:series:decide}
\begin{enumerate}
  \item Does $u_n \to 0$? If not, divergence, stop.
  \item Nonnegative terms: seek an equivalent of $u_n$ (expansions,
        \cref{ch:b1:taylor}!), compare with Riemann or geometric
        scales; factorials and powers call for the ratio test;
        decreasing $f(n)$ calls for integral comparison.
  \item Signs vary: try absolute convergence first; if it fails, the
        alternating test (check \emph{decreasing} carefully); beyond
        that, second-year tools.
\end{enumerate}
\end{method}

\section{Exercises}

\begin{exercise}[$\star$]\label{exo:b1:series:1}
Nature (and sum, when telescoping) of:
\[
\sum_{n\geq1} \frac{1}{n(n+1)},
\qquad
\sum_{n\geq2} \ln\Bigl(1 - \frac{1}{n^2}\Bigr),
\qquad
\sum_{n\geq0} \frac{3^n + 4^n}{5^n} .
\]
\end{exercise}

\begin{exercise}[$\star$]\label{exo:b1:series:2}
Nature of: $\;\sum \dfrac{n^2}{2^n}$; $\;\sum \dfrac{n!}{n^n}$;
$\;\sum \dfrac{2^n\,n!}{n^n}$; $\;\sum \dfrac{3^n\,n!}{n^n}$.
\emph{(Ratio test; recall $\bigl(1 + \frac1n\bigr)^n \to \eu$.)}
\end{exercise}

\begin{exercise}[$\star$]\label{exo:b1:series:3}
Nature of: $\;\sum \sin\dfrac{1}{n^2}$; $\;\sum
\Bigl(1 - \cos\dfrac1n\Bigr)$; $\;\sum \dfrac{1}{\sqrt{n(n+1)}}$;
$\;\sum \dfrac{\ln n}{n^2}$ \emph{(compare with $n^{-3/2}$)}.
\end{exercise}

\begin{exercise}[$\star$]\label{exo:b1:series:4}
Prove that $\sum_{n\geq1} \frac{1}{n^2}$ converges with sum $\leq 2$,
using $\frac{1}{n^2} \leq \frac{1}{n(n-1)}$ for $n \geq 2$ and a
telescoping bound.
\end{exercise}

\begin{exercise}[$\star\star$]\label{exo:b1:series:5}
Nature of $\;\sum \dfrac{(-1)^n}{\sqrt n}$, of $\;\sum
\dfrac{(-1)^n}{n + (-1)^n}$ \emph{(expand: the alternating test does
not apply directly --- why?)}, and of $\;\sum
\sin\bigl(\pi\sqrt{n^2+1}\,\bigr)$ \emph{(reduce modulo $\pi$:
$\sqrt{n^2+1} = n + \frac{1}{2n} + O(n^{-3})$)}.
\end{exercise}

\begin{exercise}[$\star\star$]\label{exo:b1:series:6}
For which $\alpha > 0$ does $\sum_{n \geq 2}
\dfrac{1}{n (\ln n)^{\alpha}}$ converge? \emph{(Integral comparison;
substitute $u = \ln t$.)}
\end{exercise}

\begin{exercise}[$\star\star$]\label{exo:b1:series:7}
Let $u_n = \dfrac{1}{n} - \ln\Bigl(1 + \dfrac1n\Bigr)$. Prove that
$0 \leq u_n \leq \dfrac{1}{2n^2}$, that $\sum u_n$ converges, and
deduce the existence of Euler's constant\index{Euler's constant}:
\[
\gamma = \lim_{N \to \infty} \Bigl( \sum_{n=1}^{N} \frac 1n - \ln N
\Bigr) .
\]
\end{exercise}

\begin{exercise}[$\star\star$]\label{exo:b1:series:8}
Compute the sums
\[
\sum_{n=1}^{\infty} \frac{1}{n(n+2)}
\qquad\text{and}\qquad
\sum_{n=0}^{\infty} \frac{n}{2^n} .
\]
\emph{(For the first: partial fractions. For the second: compute
$\sum_{n=1}^{N} n x^{n-1}$ in closed form and let $N \to \infty$ at
$x = \frac12$.)}
\end{exercise}

\begin{exercise}[$\star\star\star$]\label{exo:b1:series:9}
(Cauchy condensation) Let $(u_n)$ be nonnegative and decreasing.
Prove that
\[
\sum_{n \geq 1} u_n \text{ converges}
\iff
\sum_{k \geq 0} 2^k\, u_{2^k} \text{ converges},
\]
by comparing packets of terms between consecutive powers of $2$.
Recover from it the Riemann criterion and
\cref{exo:b1:series:6}.
\end{exercise}

\begin{exercise}[$\star\star\star$]\label{exo:b1:series:10}
Using the integral identity of \cref{ex:b1:series:ln2} adapted to
$\frac{1}{1+t^2}$, prove Leibniz's formula
\[
\frac{\pi}{4} = \sum_{n=0}^{\infty} \frac{(-1)^n}{2n+1}
= 1 - \frac13 + \frac15 - \frac17 + \cdots
\]
with the error bound $\abs{R_N} \leq \frac{1}{2N+3}$.
\end{exercise}
